\documentclass[Screen16to9,17pt]{foils}
\usepackage{zencurity-slides}

\externaldocument{basic-tcpip-and-security-exercises}
\selectlanguage{english}

\begin{document}

% Netværk og TCP/IP for begyndere – Workshop II


% En hands-on workshop med fokus på netværks-sikkerhed for begyndere

% Få styr på fundamentet bag al internettrafik – én gang for alle
% Uanset om du arbejder med udvikling, drift, support eller sikkerhed, så bygger alt på ét grundlag: netværket. TCP/IP
% er sproget, som får computere, servere og tjenester til at tale sammen, men hvad betyder det egentlig i praksis? På
% dette kursus får du en klar og praktisk introduktion til, hvordan netværk hænger sammen. Vi pakker de tekniske
% begreber ud, viser hvordan data bevæger sig på tværs af kabler og routere. Og du lærer selv at opsætte et netværk fra
% bunden.
%
% Gennem korte oplæg og praktiske øvelser lærer du at:

%   *   forstå og arbejde med IP filtrering og basale firewalls
%   *   kende forskel på og bruge router og firewall
%   *   konfigurere en moderne router med firewall via webinterface
%   *   sætte Wi-Fi op med anbefalet sikkerhed (WPA2/WPA3)
%   *   se, hvordan routing fungerer på internettet
%   *   arbejde hands-on med netværk i praksis, både fælles og egne


% Du bliver introduceret til de vigtigste grundlæggende protokoller:
% * TCP Transmission Control Protocol - SYN pakker, 5-tupel begrebet, retning af trafik
% * UDP User Datagram Protocol - problemer med sikkerhed
% * DNS Domain Name System - en smule om DNS sikkerhed
% *   VLAN 802.1q – Virtuelle LAN og Zoner
%
% Kurset ser på IPv4 og IPv6 side om side, så du forstår både paralleller og forskelle.
%
% Efter kurset vil du kunne:
%
%   *   opsætte og fejlfinde et lille netværk hjemme eller i en mindre virksomhed
%   *   forstå hvordan internettet “taler sammen”
%   *   læse netværkstrafik med fx Wireshark og Tcpdump
%   *   arbejde mere sikkert og effektivt i dit daglige it-arbejde
%
% Praktisk info
% Medbring en computer (uanset styresystem) og en browser – mere skal der ikke til.
% Ingen forudgående kendskab til netværk eller TCP/IP kræves.
%
% Anbefalet bog men det er ikke et krav at købe denne:
% Practical Packet Analysis – Using Wireshark to Solve Real-World Network Problems
% 3rd edition (2017), Chris Sanders. ISBN: 9781593278021
% Nøgleord: TCP/IP, TCP, UDP, DHCP, DNS, ICMP, TLS, Wi-Fi, WPA2, VLAN



\mytitlepage
{Networking and TCP/IP for Beginners\\Workshop II: Security}
{PROSA November 2026}


\hlkprofiluk

\slide{Goals for today}

\hlkimage{6cm}{bornhack-camp-2024.jpg}

\begin{list2}
%\item How Internet routing works
\item Get to know the difference between a router and a firewall
%\item Understand the filtering available in routers and firewalls
\item How to configure a firewall with a web interface -- OPNsense as the example
\item Discuss Wi-Fi Security -- WPA2/WPA3
\item Work hands-on with your systems and IP protocols
\item Connecting to networks bring security risks, sorry -- so we will introduce a little security
\end{list2}

Photo is NWWC camp at BornHack 2024, next is BornHack July 21. to 28. 2027 \link{https://bornhack.dk}


\slide{Time schedule}

\hlkimage{4cm}{kame-noanime-small.png}
\begin{list2}
\item 17:00 - 18:15 Introduction routing and filtering basics -- with my networks
\item 30min Break, sandwich, talk or start exercises, play with my networks, see devices
\item 18:45 - 19:45 More Security and Wi-Fi security
\item 15min small break
\item 20:00 - 21:00 Digging into some examples of programs that help
\end{list2}

During the presentation, in breaks etc. please do connect to the network, play with TCP/IP, switches and routers.

\slide{Session 0: Introduction routing and filtering basics}


\slide{Exercises are completely optional}

We will use a combination of your systems, my networking hardware and my systems.\\
{\color{red} \bf There might be live sniffing done on traffic!}

We try to mimic what you would do in your own networks during the exercises.

\begin{list2}
\item Try ping and traceroute
\item See your own IP settings
\item Borrow a USB Ethernet and connect to a switch or router
\item Borrow a router home -- only brought a few!
\end{list2}

\slide{Linux is a toolbox}

I will use Linux and participants are welcome to use any operating system of their choice.

Note: even though I talk a lot about Unix, you can definitely run a lot of tools on Windows and Mac OS X.

\begin{list2}
\item Ping and traceroute - they are on every OS
\item USB Ethernet work on most operating system
\item USB Wi-Fi cards are fun, but some takes a little work to get going
\item Basic tools like the Nmap suite is available for Microsoft Windows, Mac OS X, Linux and many others
\end{list2}

Command line tools are sometimes used in the slides, as they only show text where a GUI screenshot can be cluttered with a lot of information, feel free to find GUI tools and web sites with same functionality



\slide{Prerequisites}

\textbf{\LARGE If you are interested in TCP/IP you are welcome}

\vskip 1cm
\begin{list1}
\item It is recommended to use virtual machines for the exercises
\item Network security and most internet related security work has the following requirements:
\begin{list2}
\item Network experience
\item TCP/IP principles - often in more detail than a common user
\item Programming is an advantage, for automating things
\item Some Linux and Unix knowledge is in my opinion a {\bf necessary skill} for infosec work -- too many new tools to ignore, and lots found at sites like Github and Open Source written for Linux
\end{list2}
\end{list1}

\slide{Course Materials}

\begin{list2}
\item This material is in multiple parts:

\item Slide show - presentation - this file
\item Exercises - PDF which is used for this and other workshops
\item Additional resources from the internet are linked throughout
\item Wikipedia has a LOT of nice pages about IP protocols, for example:
\end{list2}

\begin{quote}\small
Transport Layer Security (TLS) is a cryptographic protocol designed to provide communications security over a computer network. The protocol is widely used in applications such as email, instant messaging, and voice over IP, but its use in securing HTTPS remains the most publicly visible.
\end{quote}
Source: \url{https://en.wikipedia.org/wiki/Transport_Layer_Security}

\slide{Book: Practical Packet Analysis (PPA)}


\hlkimage{6cm}{PracticalPacketAnalysis3E_cover.png}

\emph{Practical Packet Analysis,
Using Wireshark to Solve Real-World Network Problems}
by Chris Sanders, 3rd Edition
April 2017, 368 pp.
ISBN-13:
978-1-59327-802-1
\link{https://nostarch.com/packetanalysis3}

I recommend this book for people new to networking, it has been in HumbleBundle book bundles multiple times


\slide{Book: Applied Network Security Monitoring (ANSM)}

\hlkimage{5cm}{ansm-book.png}

\emph{Applied Network Security Monitoring: Collection, Detection, and Analysis}
1st Edition

Chris Sanders, Jason Smith
eBook ISBN: 9780124172166
Paperback ISBN: 9780124172081 496 pp.
Imprint: Syngress, December 2013

{\footnotesize\link{https://www.elsevier.com/books/applied-network-security-monitoring/unknown/978-0-12-417208-1}}




\slide{Internet Today}

\begin{quote}
\textbf{\large Network security}

Network security is an umbrella term to describe \textbf{security controls, policies, processes and practices} adopted to \textbf{prevent, detect and monitor unauthorized access, misuse, modification, or denial} of a \textbf{computer network and network-accessible resources}.[1] Network security involves the authorization of access to data in a network, which is controlled by the network administrator.

\end{quote}
Source: \url{https://en.wikipedia.org/wiki/Network_security}

Note: protect BOTH network and network-accessible resources

\slide{Course Network}
.
\hlkrightpic{65mm}{-1cm}{sample-network.png}

\begin{list1}
\item I have a course network with me which has the following information:
\begin{list2}
\item Routers -- computer with OpenBSD and portable routers
\item Switches Juniper and TP-Link
\item wireless access-points Unifi AP
% \item IPv4 addresses: 151.216.32.0/21 total 2048
% \item IPv6 addresses route6: 2001:678:9ec::/48
% \item Autonomous system number: AS208647 BornHack\\
\url{https://en.wikipedia.org/wiki/Autonomous_system_(Internet)}

\end{list2}
\end{list1}

\centerline{\Large You are encouraged to use the network}


\slide{Internet Today}

\hlkimage{10cm}{images/server-client.pdf}

\begin{list1}
\item Clients and servers, roots in the academic world
\item Protocols are old, some more than 40 years
\item Very few protocols were encrypted, today a lot has switched to HTTPS and TLS
\end{list1}

\slide{What the internet is actually made up of}


\hlkimage{13cm}{how-the-internet-works.jpg}
{\small Source: \emph{How the Internet Really Works
An Illustrated Guide to protocols, privacy, censorship, and governance} by catnip and article 19}\\
\url{https://catnip.article19.org/} {\color{red}\faHeart}
fun and educational book



\slide{Protection, building secure and robust networks}

\hlkimage{16cm}{sample-ip-network.pdf}


\begin{list2}
\item Network traffic is sent using network protocols with fields, values and destinations
\item Can we inspect this and decide to pass or block traffic, certainly! Does it help, sure!
\end{list2}



\slide{Network Protocol Knowledge Needed for Network Security}

To work with network security the following protocols are the bare minimum to know about.

\begin{list2}
\item ARP Address Resolution Protocol for IPv4
\item NDP Neighbor Discovery Protocol for IPv6
\item IPv4 \& IPv6 -- the basic packet fields source, destination,
\item ICMPv4 \& ICMPv6 Internet Control Message Protocol
\item UDP User Datagram Protocol
\item TCP Transmission Control Protocol
\item DHCP Dynamic Host Configuration Protocol
\item DNS Domain Name System
\end{list2}

These protocols are part of the Internet Protocol suite, or TCP/IP for short. The canonical document describing this is from 1981 RFC-0791 \emph{Internet Protocol}. The protocols were deployed on the internet around 1983.


\slide{Common Address Space}

\vskip 2 cm
\hlkimage{13cm}{IP-address.pdf}

\begin{list2}
\item Internet is defined by the address space
\item IPv4 based on 32-bit addresses, example dotted decimal format 10.0.0.1
\item IPv6 very similar to IPv4 without NAT, 128-bit addresses in hex ::1, 2a06:d380:0:101::80
\end{list2}


\slide{How to use the Internet Protocols (IP)}

Names are used by humans
\begin{center}\hlkbig
www.kramse.org

hlk@kramse.org
\end{center}

Computers use the addresses

\begin{alltt}
www     IN      A       185.129.63.130
        IN      AAAA    2a06:d380:0:102::80
mail    IN      A       217.157.63.115
        IN      AAAA    2a06:d380:0:102::25
\end{alltt}


\slide{IP addresses as bits}

\begin{alltt}
IP-address: 127.0.0.1
Binary:	01111111000000000000000000000001
\end{alltt}

\begin{list1}
\item Computers use bits \item Binary:	01111111 . 00000000 . 00000000 . 00000001
\item 127.0.0.0/8 = subnetmask with 8 bits is the special local host prefix \textbf{127}.0.0.0/8
\item 10.0.12.0/24 has prefix \textbf{10.0.12}.0/24 - 24 subnet bits
\vskip 5mm
\item IPv6 is even easier as most subnets are /64 -- so the first half is the subnet, the rest is the host
\item 2a06:d380:0:102::25/64 becomes subnets \textbf{2a06:d380:0000:0102::/64} and host part is just the end \textbf{25}
\end{list1}
\slide{RFC-1918 Private Networks}


Address Allocation for Private Internets RFC-1918
\begin{list2}
\item 10.0.0.0    -  10.255.255.255  (10/8 prefix)
\item 172.16.0.0  -  172.31.255.255  (172.16/12 prefix)
\item 192.168.0.0 -  192.168.255.255 (192.168/16 prefix)
\end{list2}

Can be used by anyone in their own local area network (LAN)

A typical \emph{home network} would use 192.168.0.0/24 with 192.168.0.1 being the router


\slide{Your own private IPv6 prefix!}

\begin{alltt}\center\textbf{
fd21:47a9:5f29::/48}\\
This is a randomly generated 48-bit unique local IPv6 prefix
as defined by RFC 4193. It can be used for local IPv6
networking. A new random prefix will be generated every time
the regenerate button is pressed, using JavaScript.
\end{alltt}
Source: \url{https://www.unique-local-ipv6.com/}

\begin{list2}
\item RFC 4193 \emph{Unique Local IPv6 Unicast Addresses}
\item Go ahead try it
\item Only usable within your own networks -- like RFC1918
\end{list2}



\slide{Documentation Prefix, IPv6 updates etc.}

Even documentation has its own prefix, RFC5737:
\begin{alltt}
The blocks 192.0.2.0/24 (TEST-NET-1), 198.51.100.0/24 (TEST-NET-2),
and 203.0.113.0/24 (TEST-NET-3) are provided for use in
documentation.
\end{alltt}

IPv6 listed in RFC3849 \verb+2001:DB8::/32+

See RFC3330 \emph{Special-Use IPv4 Addresses} which is updated by
RFC6890 \emph{Special-Purpose IP Address Registries} which in turn is updated by RFC8190\\
Use the web version of RFCs to surf back and forth \link{https://www.rfc-editor.org/rfc/rfc8190}


\slide{A switch}

\hlkimage{10cm}{switch-1.pdf}

\begin{list1}
\item Today we use switches, Don't buy a hub, not even for experimenting or sniffing
\item A switch can receive and send data on multiple ports at the same time
\item Performance only limited by the backplane and switching chips
\item Can also often route with the same speed and mirror packets
\end{list1}

\slide{Virtual LAN (VLAN)}

\hlkimage{8cm}{vlan-portbased.pdf}

\begin{list1}
\item Managed switches often allow splitting into zones called virtual LANs
\item Most simple version is port based
\item Like putting ports 1-4 into one LAN and remaining in another LAN
\item Packets must traverse a router or firewall to cross between VLANs
\end{list1}

\slide{Virtual LAN (VLAN) IEEE 802.1q}

\hlkimage{16cm}{vlan-8021q.pdf}

\begin{list1}
\item Using IEEE 802.1q  VLAN tagging on Ethernet frames
\item Virtual LAN, to pass from one to another, must use a router/firewall
\item Allows separation/segmentation and protects traffic from many security issues
\end{list1}


\slide{Network Segmentation: Which VLANs to create}

Using VLANs to segment networks we can extend it to the rest of the network, and an example of isolated segments could be:
\begin{list2}
\item Guest network for cabled and Wi-Fi clients
\item Client networks can be isolated per department, floor or building
\item Server networks for multiple purposes – development, staging, testing, production, manufacturing, …
\item Internet servers in a demilitarized zone (DMZ) which would make it less likely that compromised internet
server would affect the rest of the network
\item Dedicated printer network
\item \textbf{Management network for virtualisation, one for network management, one for storage}
\end{list2}


\slide{A modern small router includes Wireless}

\hlkimage{7cm}{sft1200_1.png}

\begin{list2}
\item Opal (GL-SFT1200) is a pocket-sized travel router supporting 1200Mbps wireless,
Max. 300 Mbps (2.4GHz) + 867 Mbps (5GHz) Fast Wi-Fi Speeds, Powerful CPU with 3 x Gigabit Ports,
Excellent Security with VPN -- OpenVPN \& WireGuard, IPv6
\end{list2}
Source: \url{https://www.gl-inet.com/products/gl-sft1200/}


\slide{Addressing examples here}

\begin{list2}
\item Home networks and adding guest and home office VLANs
\item Small company
\item Bigger companies
\item Hosting setup
\end{list2}

Lets work through some example using RFC1918 and documentation prefixes 192.0.2.0/24, 198.51.100.0/24, 203.0.113.0/24 and 2001:db8::/32 instead of real public addresses

Most IPv4 private networks would use NAT to access the internet!



\slide{NAT Network Address Translation}
\hlkimage{16cm}{nat-1.pdf}


\vskip 2 cm
\begin{list2}
\item NAT is used for connecting private networks to the Internet
\item NAT gateway replaces source address and forwards packets
\item A quick and dirty fix that keeps messing up networks and protocols
\item The NAT router/firewall has state tables
\end{list2}

  \slide{Home network: guest and home office VLANs}

  ISP gives one public IPv4 address (NAT) and IPv6 \verb+2001:db8:1200::/56+

  \begin{alltt}\small
  \textbf{VLAN Name         IPv4 (NAT)        IPv6}
  10   Home         192.168.10.0/24   2001:db8:1200:10::/64
  20   Guest Wi-Fi  192.168.20.0/24   2001:db8:1200:20::/64
  30   Home office  192.168.30.0/24   2001:db8:1200:30::/64
  40   IoT and TV   192.168.40.0/24   2001:db8:1200:40::/64
  \end{alltt}

  \begin{list2}
  \item Let the VLAN ID match the third octet and the IPv6 subnet ID -- easy to remember
  \item An ISP would give you /48 or /56, a /56 holds 256 /64 subnets, so there is room for many more VLANs
  \item Firewall rule: Guest and IoT may reach the internet, but not Home or Home office
  \end{list2}

\slide{Small company: clients, servers and a DMZ}

ISP gives external public IP \verb+203.0.113.0/29+ and \verb+2001:db8:100::/48+, internal IPv4 uses \verb+10.10.0.0/16+

\begin{alltt}\small
\textbf{VLAN Name          IPv4              IPv6}
10   Clients       10.10.10.0/24     2001:db8:100:10::/64
11   Admin_clients 10.10.11.0/24     2001:db8:100:11::/64
20   Servers       10.10.20.0/24     2001:db8:100:20::/64
30   Printers      10.10.30.0/24     2001:db8:100:30::/64
40   Guest         10.10.40.0/24     2001:db8:100:40::/64
50   DMZ           203.0.113.0/29    2001:db8:100:50::/64
99   Management    10.10.99.0/24     2001:db8:100:99::/64
\end{alltt}

\begin{list2}
\item The /29 has 8 addresses: first and last not used network, broadcast
\item Their router uses 1 so only about 5 available for customer
\item Clients, printers and guests use NAT for IPv4 when going to the internet
\item Management only reachable from admin workstations VLAN 11
\item Hint: do not use the first networks from 10.0.0.0/16 as you may want to do Site-2-Site VPN later with others
\end{list2}

  \slide{Bigger company: multiple sites}

  Plan \verb+10.0.0.0/8+ and \verb+2001:db8:a0::/44+, one /16 and one /48 per site

  \begin{alltt}\small
  \textbf{VLAN Site    Name     IPv4             IPv6}
  10   HQ      Clients  10.1.10.0/24     2001:db8:a1:10::/64
  20   HQ      Servers  10.1.20.0/24     2001:db8:a1:20::/64
  99   HQ      Mgmt     10.1.99.0/24     2001:db8:a1:99::/64

  10   Branch  Clients  10.2.10.0/24     2001:db8:a2:10::/64
  20   Branch  Servers  10.2.20.0/24     2001:db8:a2:20::/64

  -    WAN     Links    10.255.0.0/24    2001:db8:af::/48
  \end{alltt}

  \begin{list2}
  \item Same VLAN ID means same function at every site -- one firewall policy template
  \item Each site is announced as one /16 and one /48 -- small routing tables
  \item Avoid 10.0.0.0/24 and 192.168.0.0/24 internally, home networks use them and VPN users get conflicts
  \end{list2}


\slide{Bridges and routers}

\hlkimage{17cm}{wan-network.pdf}


  \slide{Hosting setup: a VLAN per customer}

  Hosting provider has \verb+198.51.100.0/24+ and \verb+2001:db8:5000::/48+

  \begin{alltt}\small
  \textbf{VLAN Name         IPv4                IPv6}
  101  Customer A   198.51.100.0/28     2001:db8:5000:100::/56
  102  Customer B   198.51.100.16/28    2001:db8:5000:200::/56
  200  Shared web   198.51.100.64/26    2001:db8:5000:300::/56
  900  Management   10.90.0.0/24        2001:db8:5000:900::/56
  901  Storage      10.91.0.0/24        2001:db8:5000:f00::/56
  \end{alltt}

  \begin{list2}
  \item IPv4 is scarce: a /28 gives 14 usable addresses, the gateway takes one of them
  \item IPv6 is plentiful: every customer gets a /56 -- 256 /64 subnets each
  \item Point-to-point links use /31 (RFC 3021) and /127 (RFC 6164)
  \item Management and storage must never be reachable from customer VLANs
  \end{list2}


\slide{Hosting and internet providers}

\hlkimage{16cm}{network-bgp-asn.png}

\begin{list2}
\item BGP networks are used for all of the Internet
\item New standards like Resource Public Key Infrastructure (RPKI) are underway
\item Try RIPE BGPlay \url{https://stat.ripe.net/special/bgplay#bgplay_fetch.resource=185.129.60.1}
\end{list2}


\slide{Routing Basics: Prefixes, AS and BGP}

\begin{list2}
\item {\bf Prefix} -- a block of addresses, like 192.0.2.0/24 or 2001:db8:1234::/48
\item {\bf Autonomous System (AS)} -- a network under one routing policy, identified by an AS number (ASN) registered at an RIR such as RIPE NCC
\item {\bf BGP} -- Border Gateway Protocol, how autonomous systems tell each other which prefixes they can reach
\item {\bf RPKI and ROA} -- signed statements of which AS may originate a prefix, prefixes without a valid ROA are filtered by many networks
\item {\bf Upstream / transit} -- the provider that announces your prefixes to the rest of the internet
\item {\bf Smallest routable prefix} -- normally /24 for IPv4 and /48 for IPv6
\item Using your own PI addresses requires an ASN, a BGP session and an upstream willing to announce them -- usually a business connection.
\item Exception: Netnørden offers BGP on private connections \link{https://netnoerden.dk/faq}
\end{list2}


\slide{DHCP Dynamic Host Configuration Protocol}

\hlkimage{10cm}{dhcp-1.pdf}

DHCP is responsible for giving you a dynamic address
\begin{list1}
\item How does a system get the address?
\item Typically in  IPv4 they will get them from DHCP
\item System sends a DHCP request, server allocates an address \emph{lease}
\item IPv6 typically uses Router Advertisement -- get prefix, configure their own IPv6 address
\end{list1}


\slide{Packets across the wire or wireless}

\hlkimage{20cm}{ethernet-frame-1.pdf}
\begin{list1}
\item Looking at data as a stream the packets are a pattern laid on top
\item Network technology defines the start and end of a frame, example Ethernet
\item From a lower level we receive a packet, example 1500-bytes from Ethernet driver
\item Operating system masks a lot of complexity
\end{list1}







\slide{IPv4 header - RFC-791 september 1981}

\begin{alltt}\footnotesize
    0                   1                   2                   3
    0 1 2 3 4 5 6 7 8 9 0 1 2 3 4 5 6 7 8 9 0 1 2 3 4 5 6 7 8 9 0 1
   +-+-+-+-+-+-+-+-+-+-+-+-+-+-+-+-+-+-+-+-+-+-+-+-+-+-+-+-+-+-+-+-+
   |Version|  IHL  |Type of Service|          Total Length         |
   +-+-+-+-+-+-+-+-+-+-+-+-+-+-+-+-+-+-+-+-+-+-+-+-+-+-+-+-+-+-+-+-+
   |         Identification        |Flags|      Fragment Offset    |
   +-+-+-+-+-+-+-+-+-+-+-+-+-+-+-+-+-+-+-+-+-+-+-+-+-+-+-+-+-+-+-+-+
   |  Time to Live |    Protocol   |         Header Checksum       |
   +-+-+-+-+-+-+-+-+-+-+-+-+-+-+-+-+-+-+-+-+-+-+-+-+-+-+-+-+-+-+-+-+
   |                       Source Address                          |
   +-+-+-+-+-+-+-+-+-+-+-+-+-+-+-+-+-+-+-+-+-+-+-+-+-+-+-+-+-+-+-+-+
   |                    Destination Address                        |
   +-+-+-+-+-+-+-+-+-+-+-+-+-+-+-+-+-+-+-+-+-+-+-+-+-+-+-+-+-+-+-+-+
   |                    Options                    |    Padding    |
   +-+-+-+-+-+-+-+-+-+-+-+-+-+-+-+-+-+-+-+-+-+-+-+-+-+-+-+-+-+-+-+-+

                    Example Internet Datagram Header
\end{alltt}
Source: \url{https://datatracker.ietf.org/doc/html/rfc791} and updated later

\slide{IPv6 header - RFC-1883 December 1995}
\vskip -1cm
\begin{alltt}\footnotesize
   +-+-+-+-+-+-+-+-+-+-+-+-+-+-+-+-+-+-+-+-+-+-+-+-+-+-+-+-+-+-+-+-+
   |Version| Traffic Class |           Flow Label                  |
   +-+-+-+-+-+-+-+-+-+-+-+-+-+-+-+-+-+-+-+-+-+-+-+-+-+-+-+-+-+-+-+-+
   |         Payload Length        |  Next Header  |   Hop Limit   |
   +-+-+-+-+-+-+-+-+-+-+-+-+-+-+-+-+-+-+-+-+-+-+-+-+-+-+-+-+-+-+-+-+
   |                                                               |
   +                                                               +
   |                                                               |
   +                         Source Address                        +
   |                                                               |
   +                                                               +
   |                                                               |
   +-+-+-+-+-+-+-+-+-+-+-+-+-+-+-+-+-+-+-+-+-+-+-+-+-+-+-+-+-+-+-+-+
   |                                                               |
   +                                                               +
   |                                                               |
   +                      Destination Address                      +
   |                                                               |
   +                                                               +
   |                                                               |
   +-+-+-+-+-+-+-+-+-+-+-+-+-+-+-+-+-+-+-+-+-+-+-+-+-+-+-+-+-+-+-+-+
\end{alltt}
%Source: \url{https://datatracker.ietf.org/doc/html/rfc1883} and obsoleted by\\
%\url{https://datatracker.ietf.org/doc/html/rfc8200}

\slide{UDP User Datagram Protocol}
\hlkimage{16cm}{udp-1.pdf}
Connectionless \url{https://en.wikipedia.org/wiki/User_Datagram_Protocol}


\slide{TCP Transmission Control Protocol}
\hlkimage{14cm}{tcp-1.pdf}

Connection-oriented \url{https://en.wikipedia.org/wiki/Transmission_Control_Protocol}


\slide{TCP three way handshake}

\hlkimage{6cm}{images/tcp-three-way.pdf}

\begin{list2}
\item Session setup is used in some protocols
\item TCP Fast Open and TLS 1.3 / QUIC 0-RTT can send data in the first packet, HTTP/2 still uses the full TCP handshake
\end{list2}


\slide{Well-Known Port Numbers}

\hlkimage{6cm}{iana1.jpg}

\begin{list1}
\item IANA maintains a list of magical numbers in TCP/IP
\item Lists of protocol numbers, port numbers etc.
\item A few notable examples:
\begin{list2}
\item Port 25/tcp Simple Mail Transfer Protocol (SMTP)
\item Port 53/udp and 53/tcp Domain Name System (DNS)
\item Port 80/tcp Hyper Text Transfer Protocol (HTTP)
\item Port 443/tcp HTTP over TLS/SSL (HTTPS)
\end{list2}
\item Source: \link{http://www.iana.org}
\end{list1}



\slide{Defense in depth}

%\hlkimage{10cm}{Bartizan.png}
\hlkimage{15cm}{medieval-clipart-5}
\centerline{Picture originally from: \url{http://karenswhimsy.com/public-domain-images}}



\slide{Back in the day: Firewalls and Internet Security, 1994}

\hlkimage{4cm}{images/cheswick-cover2e.jpg}
\begin{list2}
\item \emph{Firewalls and Internet Security: Repelling the Wily Hacker} , Second Edition 2003, William R. Cheswick, Steven M. Bellovin, and Aviel D. Rubin, {\bf 2003} \link{http://www.wilyhacker.com/}
\item The full PDF—and the full LaTeX source of the book.
Creative Commons Attribution-NonCommercial-NoDerivatives 4.0 International License.
\item How to configure firewalls often boil down to, should we allow protocol X
\item If we allow certain protocols through a firewall, we are asking for trouble
\end{list2}


\slide{Network Segmentation -- Firewalls}

\begin{quote}\small
\$ firewall\\

1. (I) {\bf An internetwork gateway that restricts data communication traffic to and from one of the connected networks} (the one said to be "inside" the firewall) and thus protects that network's system resources against threats from the other network (the one that is said to be "outside" the firewall). (See: guard, security gateway.)

2. (O) {\bf A device or system that controls the flow of traffic between networks using differing security postures.} Wack, J. et al (NIST), "Guidelines on Firewalls and Firewall Policy", Special Publication 800-41, January 2002.

Tutorial: A firewall typically protects a smaller, secure network (such as a corporate LAN, or even just one host) from a larger network (such as the Internet). The firewall is installed at the point where the networks connect, and the firewall applies policy rules to control traffic that flows in and out of the protected network.
\end{quote}
{\footnotesize Source: RFC4949 \emph{Internet Security Glossary, Version 2}\\
\link{https://datatracker.ietf.org/doc/html/rfc4949} 2007}

\slide{Continued}
\begin{quote}\small
{\bf A firewall is not always a single computer.} For example, a firewall may consist of a pair of filtering routers and one or more proxy servers running on one or more bastion hosts, all connected to a small, dedicated LAN (see: buffer zone) between the two routers. The external router blocks attacks that use IP to break security (IP address spoofing, source routing, packet fragments), while proxy servers block attacks that would exploit a vulnerability in a higher-layer protocol or service. The internal router blocks traffic from leaving the protected network except through the proxy servers. The difficult part is defining criteria by which packets are denied passage through the firewall, because a firewall not only needs to keep unauthorized traffic (i.e., intruders) out, but usually also needs to let authorized traffic pass both in and out.
\end{quote}
{\footnotesize Source: RFC4949 \emph{Internet Security Glossary, Version 2}\\
\link{https://datatracker.ietf.org/doc/html/rfc4949} 2007}

\begin{list2}
\item Network layer, packet filters, application level, stateless, stateful
\item Firewalls are by design a choke point, natural place to do network security monitoring!
\item Older but still interesting Cheswick chapter 2 PDF
\emph{A Security Review of Protocols:
Lower Layers}\\
\link{http://www.wilyhacker.com/}
\end{list2}


\slide{What is a packet filter}

A packet filter looks at the header of each individual packet and decides: pass or block
\begin{list2}
\item Matches on layer 3 and 4 fields: source and destination IP, protocol (TCP, UDP, ICMP), ports, TCP flags, interface and direction
\item {\bf Stateless} -- each individual packet is judged alone. On routers and switches this is an Access Control List (ACL), fast and often done in hardware
\item Stateless rules must also allow the return traffic, for example TCP packets with ACK set
\item {\bf Stateful} -- keeps a state table of connections, return packets for a known connection are allowed automatically
\item Both only see layer 3 and 4, not the content of HTTP, DNS etc. -- that requires application proxies or a Web Application Firewall (WAF)
\item Firewalls are by design a choke point, natural place to do network security monitoring!
\end{list2}

They are all firewalls -- or firewall devices!

\slide{Access Control Lists (ACL)}

Stateless firewall filter throw stuff away

\begin{alltt}\footnotesize
hlk@MX-CPH-02> show configuration firewall filter all | no-more
/* This is a sample, maybe use BGP flowspec and/or RTBH */
term edgeblocker \{
    from \{
        source-address \{
            84.180.xxx.173/32;
...
            87.245.xxx.171/32;
        \}
        destination-address \{
            91.102.xx.xx/28; \}
        protocol [ tcp udp icmp ]; \}
    then discard;
\}
\end{alltt}
Hint: can also leave out protocol and then it will match all protocols

\slide{Strict Filtering for Segments -- still stateless!}

\begin{alltt}\footnotesize
term some-server-allow \{
    from \{
        destination-address \{
            185.129.xx.0/xx;
        \}
        protocol tcp;
        destination-port [ 80 443 ];
    \} then accept;
\}
term some-server-block-unneeded \{
    from \{
        destination-address \{
            185.129.xx.0/xx; \}
        protocol-except icmp;  \}
    then \{ count some-server-block; discard;
    \}
\}
\end{alltt}

Wut - no UDP, yes UDP service is not used on these servers

\slide{Control-Plane Access Control Lists (CP-ACL)}

Management of network devices can be done with Access Control lists, other systems may be put into a management VLAN

An example from a Cisco router using secure shell is shown below with a simple standard Access Control List (ACL) named 22 and then referenced for the virtual terminal (vty) secure shell (ssh):

\begin{alltt}\scriptsize
ip access-list standard 22
10 permit 192.0.2.2
20 permit 172.16.22.10
30 permit 192.168.0.10
40 permit 203.0.113.10
line vty 0 4
  access-class 22 in
  transport input ssh
\end{alltt}

This will reduce the likelihood that an attacker can gain access to the administration using leaked username and passwords or because of vulnerabilities in the software.

\slide{Creating an Access Control List (ACL)}

\begin{alltt}\small
 (config)#ipv6 access-list RA-GUARD
 (config-ipv6-acl)#sequence 3 deny icmp any any router-advertisement
 (config-ipv6-acl)#sequence 6 permit ipv6 any any
 (config-ipv6-acl)#exit
 (config)#interface FastEthernet0/5
 (config-if)#ipv6 traffic-filter RA-GUARD in
\end{alltt}
Source: example copied from RIPE NCC IPv6 Security Training materials:\\
\link{https://www.ripe.net/support/training/material/ipv6-security/ipv6security-slides.pdf}

\begin{list2}
\item Best practice, and not that hard to do -- Layer 2 protection
\item ACL, filtering and firewalling will create longer lasting protection
\item Paired with a nice address plan you can easily put restrictions on traffic flow, without hurting functionality or the business
\item Does ANY client in ANY office NEEEEEED to connect to ANY UPS, Virtualisation and printer across the world ...
\end{list2}



\slide{Firewalls}

\hlkimage{8cm}{opal-luci-firewall-easy.png}

\begin{list2}
\item Firewalls are a great project to look into, there are so many \smiley
\item I recommend looking into open source firewalls like OPNsense, UFW and OpenBSD and many routers also include firewalls. You can often find images from commercial vendors, that you can try out.
\item Picture is from GL-Inet router Opal with OpenWRT/LuCI which is router software for smaller routers
\end{list2}


\slide{Modern Firewall Infrastructures}


\centerline{\hlkbig A firewall {\color{security6blue}blocks traffic} on a network}

\vskip 1 cm
\pause

\centerline{\hlkbig A firewall {\color{red}allows traffic} on a network}
{\small The interesting part is typically what it allows!}

\begin{list1}
\item A firewall infrastructure must:
\begin{list2}
\item Prevent attackers from entering
\item Prevent data exfiltration
\item Prevent worms, malware, virus from spreading in networks
\item Be part of an overall solution with ISP, routers, other firewalls, switched infrastructures,
  intrusion detection systems and the rest of the infrastructure
\item ...
\end{list2}
\end{list1}

\vskip 5mm
\centerline{Difficult -- and requires design and secure operations}


\slide{Modern Firewalls}

\begin{list1}
\item Basically some filtering between networks or network segments
\item Typically they contain:
  \begin{list2}
   \item Some interface, maybe web interface, often command line interface
\item TCP/IP filtering options -- packets flowing in and out, direction, protocol, ports etc.
\item Should be able to handle both IPv6 and legacy IPv4
\item Often they have predefined rules for common use-cases. Is this really a good thing if you can easily configure a bad protocol like Server Message Block to and from the Internet?
\item Most legacy setups use Network Address Translation (NAT) -- NAT is a kludge and bad!
\item Most platforms have extra network related features DHCP servers, DNS caching servers etc.
\end{list2}
\item The firewall devices are mostly allowing some {\bf stateful filtering} which are much easier to configure than a pure network packet filter
\end{list1}

Goal is to implement rules -- a security policy for isolation and data flow

\slide{Sample rules from OpenBSD PF}

\begin{alltt}\tiny
# hosts and networks
router="192.0.2.1"
webserver="192.0.2.80"
homenet="{ 198.51.100.0/24, 203.0.113.0/24 }"
wlan="198.51.100.0/24"
wireless=wi0
set skip lo0
# things not used
spoofed="{ 127.0.0.0/8, 10.0.0.0/8, 172.16.0.0/12, 192.168.0.0/16, 255.255.255.255/32 }"
{\bf
# default block anything
block in all }
# egress and ingress filtering - disallow spoofing, and drop spoofed
block in quick from $spoofed to any
block out quick from any to $spoofed

pass in on $wireless proto tcp from \{ $wlan $homenet \} to any port = 22
pass in on $wireless proto tcp from any to $webserver port = 80

pass out
\end{alltt}


\slide{Example Firewall Products}
\begin{list2}
\item Cisco Firewall \link{http://www.cisco.com}
\item Juniper SRX \link{http://www.juniper.net}
\item Palo Alto \link{https://www.paloaltonetworks.com/}
\item Fortinet \link{https://www.fortinet.com/}
\item Checkpoint Firewall-1 \link{http://www.checkpoint.com}
\item Multiple others exist
\end{list2}

Those listed are the most popular commercial ones I see in Denmark


\slide{Open source based firewalls}
\begin{list2}
\item Linux firewalls IP tables, use command line tool ufw Uncomplicated Firewall!
\item Firewall GUIs on top of Linux -- lots! Some are also available as commercial ones
\item OPNsense a great modern firewall \link{https://opnsense.org/}
\item OpenBSD PF
\link{http://www.openbsd.org}
\item FreeBSD IPFW and IPFW2 \link{http://www.freebsd.org}
\item Mac OS X uses OpenBSD PF
\item FreeBSD has an older version of the OpenBSD PF, should really be renamed now
\end{list2}


\slide{Firewall OPNsense GUI based and easy to install}

\hlkimage{12cm}{images/screenshots_OPNsense-1024x518.png}

\begin{list2}
\item OPNsense \link{https://opnsense.org/}
\item Firewall built on FreeBSD with web interface
\item Originally thoughts from m0n0wall and later \link{https://www.pfsense.org/}\\
\item Danish companies have been using these for many years now
\end{list2}


\slide{Linux Uncomplicated Firewall (UFW)}

Debian / Ubuntu:
\begin{alltt}\scriptsize
root@debian01:~# apt install ufw
...
root@debian01:~# ufw allow 22/tcp
Rules updated
Rules updated (v6)
root@debian01:~# ufw enable
Command may disrupt existing ssh connections. Proceed with operation (y|n)? y
Firewall is active and enabled on system startup
root@debian01:~# ufw status numbered
Status: active

     To                         Action      From
     --                         ------      ----
[ 1] 22/tcp                     ALLOW IN    Anywhere
[ 2] 22/tcp (v6)                ALLOW IN    Anywhere (v6)
\end{alltt}

\begin{list2}
\item Extremely easy to use -- I recommend and use the (Uncomplicated Firewall) UFW
\end{list2}

\slide{Session 1: More Security and Wi-Fi security}


\slide{Cilium overview}

\hlkimage{12cm}{cilium-overview.png}

\begin{quote}
Kubernetes provides Network Policies for controlling traffic going in and out of the pods. Cilium implements the Kubernetes Network Policies for L3/L4 level and extends with L7 policies for granular API-level security for common protocols such as HTTP, Kafka, gRPC, etc
\end{quote}
Source: picture and text from \link{https://cilium.io/blog/2018/09/19/kubernetes-network-policies/}

\slide{Kubernetes setup only VPN access to Nextcloud}
\hlkimage{13cm}{zencurity-k8s-netvaerk.png}
I do not use all Nextcloud features, so it is only reachable through WireGuard VPN


\slide{Firewalls are NOT Alone}

\hlkimage{15cm}{network-layers-1.png}

\centerline{Use Defense in Depth -- all layers have features}



\slide{NGINX as a Reverse Proxy}

\hlkimage{15cm}{nginx_reverse_proxy.png}
\begin{list2}
\item A reverse proxy sits in front of backend servers, accepting client requests and forwarding them internally
\item NGINX uses the \verb+proxy_pass+ directive to route traffic to upstream services by hostname or IP
\item Can provide TLS termination, load balancing, caching, and request buffering at the edge. A single system with NGINX get certificates with Let's Encrypt (or alternatives)
\end{list2}

% \link{https://en.wikipedia.org/wiki/Reverse_proxy}

\slide{Proxy servers and Web Application Firewalls (WAF)}

Filtering at higher layers is also possible
\begin{list2}
\item Web proxies for clients can help security a lot -- a centralized filter for everyone

\item Reverse proxies for web applications are called
Web Application Firewalls (WAF) -- and filter incoming web requests, and outgoing answers. Can help with attacks like SQL injection and exfiltration of data
\item Depending on your network it can replace or be combined with filtering on DNS servers, and I would prefer to filter domains with DNS
\item I would also prefer blocking large prefixes of IP destinations using routers/stateless packet filters -- maybe use BGP for distributing \emph{lists} or automation for downloading block lists
\end{list2}

Also great home lab projects!

\slide{CrowdSec -- Crowd sourced security information}

%\hlkimage{}{}

\begin{quote}
{\bf Detect and block Log4j exploitation attempts with CrowdSec}\\
If you work in Infosec, you had a very lousy weekend. And that’s because of the Log4j zero-day vulnerability (CVE-2021-44228) that was discovered. We had no choice but to roll up our sleeves to help our community before things got messier than they already were.

As a result, we have released a scenario that will help you detect and block exploitation attempts of the vulnerability. This new scenario can be directly downloaded from our Hub and installed in a blink of an eye. Check this quick video to see the plugin in action:
\end{quote}

\url{https://www.crowdsec.net/blog/detect-block-log4j-exploitation-attempts}


\begin{list2}
\item Log4j is a popular software library in the Java world
\item CrowdSec quickly provided a list of systems scanning the internet for this vulnerability\\
 \link{https://gist.github.com/blotus/f87ed46718bfdc634c9081110d243166}
\item Full disclosure they gave me a hoodie at an event \smiley
\end{list2}

\slide{CrowdSec with NGINX}

%\hlkimage{8cm}{crowdsec_nginx.png}

\begin{list2}
\item {\bf CrowdSec} is an open-source, community-driven security engine -- the llama company \link{https://crowdsec.net}
\item {\bf NGINX Bouncer} -- a Lua plugin enforcing ban decisions in real time via \verb+access_by_lua_block+,
blocking malicious IPs before requests reach the backend
\item {\bf Security Engine} -- reads NGINX logs and detects behavioural patterns: HTTP probing,
scanning, brute force, and matches against the community blocklist
\item {\bf AppSec/WAF component} -- inline virtual patching, SQLi, XSS, and known CVE protection;
install with:
\verb+cscli collections install crowdsecurity/appsec-virtual-patching+
\item {\bf Community Blocklist} -- crowd-sourced IP reputation feed, continuously updated from
millions of sensors worldwide; free tier available
\item Two complementary modes: reactive (log-based, bans offenders) and
inline (WAF, blocks the request before it is processed)
\end{list2}
Source: Docs: \link{https://docs.crowdsec.net/u/bouncers/nginx/}


\slide{The OWASP Core Rule Set for WAF}

\begin{list2}
\item {\bf OWASP Core Rule Set (CRS)} -- the standard open source rules, protects against SQL injection, XSS, remote code execution and other common attacks \link{https://coreruleset.org}
\item CRS uses {\bf anomaly scoring} -- each matching rule adds points, the request is blocked when the total passes a threshold. {\bf Paranoia levels} 1--4 -- higher levels catch more attacks, but also block more legitimate traffic, start at level 1
\item Expect tuning: false positives are normal, run in detection-only mode first and read the logs
\item {\bf ModSecurity} -- open source Web Application Firewall engine, inspects each HTTP request and response against rules. Works with Apache, NGINX and IIS, version 3 (libmodsecurity) uses a connector module for NGINX
\item Trustwave ended support in July 2024, ModSecurity is now an OWASP project\\
\link{https://owasp.org/www-project-modsecurity/}
\item {\bf Coraza} -- newer OWASP WAF engine written in Go, first stable release 2021, fully compatible with CRS v4\\
\link{https://coraza.io}

\end{list2}





\slide{First or Last match firewall?}

\hlkimage{18cm}{images/first-last-match-1.pdf}

\slide{First match - example IPFW}

\begin{alltt}
\hlksmall
00100 16389  1551541 allow ip from any to any via lo0
00200     0        0 deny log ip from any to 127.0.0.0/8
00300     0        0 check-state
...
{\bfseries
65435    36     5697 deny log ip from any to any}
65535   865    54964 allow ip from any to any
\end{alltt}

This is a stateful filtering ruleset, rule 300 checks if an existing connection is known matching the packet incoming

The deny rule in bold show that the current configuration is a \emph{best current practice} default deny

\slide{Last match - example OpenBSD PF}

\begin{alltt}\small
ext_if="ext0"
int_if="int0"
{\bf
block in}
pass out keep state

pass quick on \{ lo $int_if \}

# Allow port http=80 and port domain=53
# on the IP of the IP of the external interface ($ext_if)
pass in on $ext_if proto tcp to ($ext_if) port http keep state
pass in on $ext_if proto \{ tcp, udp \} to ($ext_if) port domain keep state
\end{alltt}

This is also a default deny, all packets incoming are marked as block, so unless they are allowed by a later rule -- thrown away



\slide{Wi-Fi Introduction and Security}



\begin{list1}
\item Wifi standards IEEE 802.11
\begin{list2}
\item 802.11b 11Mbps and 802.11g 54Mbps
\item 802.11n Wi-Fi 4, 802.11ac Wi-Fi 5
\item 802.11ax Wi-Fi 6 and Wi-Fi 6E (6GHz band)
\item 802.11be Wi-Fi 7
\item 802.11i Security enhancements Robust Security Network RSN
\end{list2}
\end{list1}

Source: \link{http://grouper.ieee.org/groups/802/11/index.html}




\slide{Wi-Fi Labs Network cards}

\hlkimage{16cm}{TL-WN722N.png}

\begin{list1}
\item Laptop or Netbook, I typically use USB wireless cards. Cheap and easy -- keep using your built-in with the host OS
\item Access Points - get a small selection for testing
\item Many resources available, books, sites, blogs, etc.
\end{list1}


\slide{IEEE 802.11 Security fast forward }

\begin{quote}
{\bf In 2001}, a group from the University of California, Berkeley presented a paper describing weaknesses in the 802.11 Wired Equivalent Privacy (WEP) security mechanism defined in the original standard; they were followed by {\bf Fluhrer, Mantin, and Shamir's} paper titled "Weaknesses in the Key Scheduling Algorithm of RC4". Not long after, Adam Stubblefield and AT\&T publicly announced the first {\bf verification of the attack}. In the attack, they were able to intercept transmissions and gain unauthorized access to wireless networks.
\end{quote}
Source: \link{http://en.wikipedia.org/wiki/IEEE_802.11}

\slide{IEEE 802.11 Security fast forward }

\begin{quote}
The IEEE set up a dedicated task group to create a replacement security solution, {\bf 802.11i} (previously this work was handled as part of a broader 802.11e effort to enhance the MAC layer). The Wi-Fi Alliance announced an {\bf interim specification called Wi-Fi Protected Access (WPA)} based on a subset of the then current IEEE 802.11i draft. These started to appear in products in {\bf mid-2003}. {\bf IEEE 802.11i (also known as WPA2)} itself was ratified in {\bf June 2004}, and uses government strength encryption in the {\bf Advanced Encryption Standard AES,} instead of RC4, which was used in WEP. The modern recommended encryption for the home/consumer space is {\bf WPA2 (AES Pre-Shared Key) and for the Enterprise space is WPA2 along with a RADIUS authentication server} (or another type of authentication server) and a strong authentication method such as EAP-TLS.
\end{quote}
Source: \link{http://en.wikipedia.org/wiki/IEEE_802.11}

\slide{IEEE 802.11 Security fast forward }

\begin{quote}
In January 2005, the IEEE set up yet another task group "w" to protect management and broadcast frames, which previously were sent unsecured. Its standard was published in 2009.[24]

In {\bf December 2011}, a security flaw was revealed that affects wireless routers with the {\bf optional Wi-Fi Protected Setup (WPS)} feature. While WPS is not a part of 802.11, {\bf the flaw allows a remote attacker to recover the WPS PIN and, with it, the router's 802.11i password in a few hours}.
\end{quote}
Source: \link{http://en.wikipedia.org/wiki/IEEE_802.11}

\vskip 1cm
\centerline{WPS WTF?! Turn it off -- use WPA PSK or WPA Enterprise}


\slide{WPA3 Security}

%\hlkimage{}{}

\begin{quote}\small
{\bf WPA3}\\
In January 2018, the Wi-Fi Alliance announced WPA3 as a replacement to WPA2.[12][13] Certification began in June 2018,[14] and WPA3 support has been mandatory for devices which bear the "Wi-Fi CERTIFIED™" logo since July 2020.[11]

The new standard uses an equivalent 192-bit cryptographic strength in WPA3-Enterprise mode[15] (AES-256 in GCM mode with SHA-384 as HMAC), and still mandates the use of CCMP-128 (AES-128 in CCM mode) as the minimum encryption algorithm in WPA3-Personal mode.

The WPA3 standard also replaces the pre-shared key (PSK) exchange with Simultaneous Authentication of Equals (SAE) exchange, a method originally introduced with IEEE 802.11s, resulting in a more secure initial key exchange in personal mode[16][17] and forward secrecy.[18] The Wi-Fi Alliance also says that WPA3 will mitigate security issues posed by weak passwords and simplify the process of setting up devices with no display interface.[2][19]

Protection of management frames as specified in the IEEE 802.11w amendment is also enforced by the WPA3 specifications.
\end{quote}
Source: \link{https://en.wikipedia.org/wiki/Wi-Fi_Protected_Access}

\begin{list2}
\item Use WPA3, or WPA2/WPA3 transition mode if you still have older devices
\end{list2}

\slide{IS WPA3 supported?}

%\hlkimage{}{}

\begin{quote}
A wireless network adapter that supports Wi-Fi 6. To see if your PC supports it, check the documentation that came with it or check the PC manufacturer’s website.
Tip: You can also check to see if your router supports Wi-Fi 6 by opening the Command Prompt, and then typing the command netsh wlan show drivers. Look next to Radio types supported and see if it includes 802.11ax.
\end{quote}
Source: \link{https://support.microsoft.com/en-us/windows/faster-and-more-secure-wi-fi-in-windows-26177a28-38ed-1a8e-7eca-66f24dc63f09}

\begin{list2}
\item Your devices must support both WPA3 in both operating system and Wi-Fi drivers!
\item Windows 10 and 11 does, so try upgrading drivers
\item Android does since Android 10 (2019)
\item Apple devices have support in recent versions
\end{list2}


\slide{WPA3 on Apple Devices}

%\hlkimage{}{}

\begin{quote}
\begin{list2}
\item iPhone 7 or later.
\item iPad 5th generation or later.
\item Apple TV 4K or later.
\item Apple Watch series 3 or later.
\item Mac computers (late 2013 or later, with 802.11ac or later)
\end{list2}
\end{quote}
Source: \link{https://support.apple.com/da-dk/guide/security/sec8a67fa93d/web}


\slide{Wi-Fi CERTIFIED WPA3™}

%\hlkimage{}{}

\begin{quote}

WPA3™ provides cutting-edge security protocols to the market. Building on the widespread success and adoption of Wi-Fi security, WPA3 adds new features to simplify Wi-Fi security, enable more robust authentication, deliver increased cryptographic strength for highly sensitive data markets, and maintain resiliency of mission critical networks. All WPA3 networks:

\begin{list2}
\item Use the latest security methods
\item Disallow outdated legacy protocols
\item Require use of Protected Management Frames (PMF)
\end{list2}
Since Wi-Fi networks differ in usage purpose and security needs, WPA3 includes additional capabilities specifically for personal and enterprise networks. Users of WPA3-Personal receive increased protections from password guessing attempts, while WPA3-Enterprise users can now take advantage of higher-grade security protocols for sensitive data networks.

WPA3 is a mandatory certification for Wi-Fi CERTIFIED™ devices.
\end{quote}
Source: \url{https://www.wi-fi.org/discover-wi-fi/security}

\slide{WPA3-Personal}

\begin{quote}
WPA3-Personal brings better protections to individual users by providing more robust password-based authentication, even when users choose passwords that fall short of typical complexity recommendations. This capability is enabled through Simultaneous Authentication of Equals (SAE). The technology is resistant to offline dictionary attacks where an adversary attempts to determine a network password by trying possible passwords without further network interaction.

\begin{list2}
\item {\bf Natural password selection}: Allows users to choose passwords that are easier to remember
\item {\bf Ease of use}: Delivers enhanced protections with no change to the way users connect to a network
\item {\bf Forward secrecy}: Protects data traffic even if a password is compromised after the data was transmitted
\end{list2}
\end{quote}

Source: \url{https://www.wi-fi.org/discover-wi-fi/security}


\slide{WPA3-Enterprise}

%\hlkimage{}{}

\begin{quote}
WPA3-Enterprise builds upon the foundation of WPA2-Enterprise with the additional requirement of using Protected Management Frames on all WPA3 connections.

\begin{list2}
\item {\bf Authentication}: multiple Extensible Authentication Protocol (EAP) methods
\item {\bf Authenticated encryption}: minimum 128-bit Advanced Encryption Standard Counter Mode with Cipher Block Chaining Message Authentication (AES-CCMP 128)
\item {\bf Key derivation and confirmation}: minimum 256-bit Hashed Message Authentication Mode (HMAC) with Secure Hash Algorithm (HMAC-SHA256)
\item {\bf Robust management frame protection}: minimum 128-bit Broadcast/Multicast Integrity Protocol Cipher-based Message Authentication Code (BIP-CMAC-128)
\end{list2}
\end{quote}

Source: \url{https://www.wi-fi.org/discover-wi-fi/security}


\slide{WPA3-Enterprise with 192-bit mode}

\begin{quote}

WPA3-Enterprise also offers an optional mode using 192-bit minimum-strength security protocols and cryptographic tools to better protect sensitive data.

\begin{list2}
\item {\bf Authentication}: Extensible Authentication Protocol – Transport Layer Security (EAP-TLS) using Elliptic Curve Diffie-Hellman (ECDH) exchange and Elliptic Curve Digital Signature Algorithm (ECDSA) using a 384-bit elliptic curve
\item {\bf Authenticated encryption}: 256-bit Galois/Counter Mode Protocol (GCMP-256)
\item {\bf Key derivation and confirmation}: 384-bit Hashed Message Authentication Mode (HMAC) with Secure Hash Algorithm (HMAC-SHA384)
\item {\bf Robust management frame protection}: 256-bit Broadcast/Multicast Integrity Protocol Galois Message Authentication Code (BIP-GMAC-256)
\end{list2}
The 192-bit security mode offered by WPA3-Enterprise ensures the right combination of cryptographic tools are used and sets a consistent baseline of security within a WPA3 network.
\end{quote}

Source: \url{https://www.wi-fi.org/discover-wi-fi/security}




\slide{2020: Dragonblood}

%\hlkimage{}{}

\begin{quote}\small
April 2019 — Modern Wi-Fi networks use WPA2 to protect transmitted data. However, because {\bf WPA2 is more than 14 years old}, the Wi-Fi Alliance recently announced the new and more secure WPA3 protocol. One of the supposed advantages of WPA3 is that, thanks to its underlying {\bf Dragonfly handshake}, it's {\bf near impossible to crack the password} of a network. Unfortunately, we found that even with WPA3, an {\bf attacker within range of a victim can still recover the password}. If the victim uses no extra protection such as HTTPS, this allows an attacker to steal sensitive information such as passwords and emails. We hope our disclosure motivates vendors to mitigate our attacks before WPA3 becomes widespread.

...

Fortunately, as a result of our research, both the Wi-Fi standard and EAP-pwd are being updated with a more secure protocol. Although this update is not backwards-compatible with current deployments of WPA3, it does prevent most of our attacks.
\end{quote}
Source: \link{https://wpa3.mathyvanhoef.com/}


\begin{list2}
\item Side-channel leaks
\item Full paper Mathy Vanhoef and Eyal Ronen. 2020. \emph{Dragonblood: Analyzing the Dragonfly Handshake of WPA3 and EAP-pwd}. In IEEE Symposium on Security \& Privacy (SP). IEEE.
 \link{https://eprint.iacr.org/2019/383}

\end{list2}



\slide{Session 2: Digging into some examples of programs that help}



\slide{Intrusion Kill Chains}

\hlkimage{13cm}{crafting-cip-kill-chain.png}

\begin{list2}
\item See also \emph{Intelligence-Driven Computer Network Defense Informed by Analysis of Adversary Campaigns and Intrusion Kill Chains}, Eric M. Hutchins , Michael J. Cloppert, Rohan M. Amin, Ph.D. Lockheed Martin Corporation\\{\footnotesize
 \link{https://www.lockheedmartin.com/content/dam/lockheed-martin/rms/documents/cyber/LM-White-Paper-Intel-Driven-Defense.pdf}}
\end{list2}




\slide{Metadata -- enrichment}

\hlkimage{10cm}{crafting-security-playbook-metadata.png}

Source: picture from Crafting the InfoSec Playbook, CIP

Metadata + Context


\slide{The Spamhaus Don't Route Or Peer Lists}

\begin{quote}
The Spamhaus Don't Route Or Peer Lists

DROP (Don't Route Or Peer) and EDROP are advisory "drop all traffic" lists, consisting of stolen 'hijacked' netblocks and netblocks controlled entirely by criminals and professional spammers. DROP and EDROP are a tiny subset of the SBL designed for use by firewalls and routing equipment.
\end{quote}

\link{http://www.spamhaus.org/drop/}


\begin{list2}
\item EDROP was merged into DROP on 10 April 2024, use only DROP and DROPv6 today
\item When your SIEM alerts you, you need tools to block and restrict
\item Recommend adding empty blocking access control lists etc. to your network infrastructure
\item Add premade blocking to your name servers, proxy servers, recursive servers
\item Recommend implementing country lists
\end{list2}



\slide{The basic tools for countering threats}

{\Large Knowledge and insight}
\begin{list2}
\item Networks have end-points and conversations on multiple layers
\item Wireshark is advanced, try right-clicking different places
\item Name resolution includes low level MAC addresses, and IP - names
\end{list2}

\begin{list2}
\item Tcpdump format, built-in to many network devices
\item Remote packet dumps, like \verb+tcpdump –i eth0 –w packets.pcap+
\item Story: tcpdump was originally written in 1988 by Van Jacobson, Sally Floyd, Vern Paxson and Steven McCanne who were, at the time, working in the Lawrence Berkeley Laboratory Network Research Group\\
 \link{https://en.wikipedia.org/wiki/Tcpdump}
\end{list2}

\vskip 5mm

\centerline{\Large Great network security comes from knowing networks!}



\slide{Wireshark - graphical network sniffer}

\hlkimage{13cm}{images/wireshark-http.png}

\centerline{Capture - Options, select a network interface}
\centerline{\link{http://www.wireshark.org}}



\slide{All attacks have signatures, some more noisy than others}

\hlkimage{14cm}{ppa-passive-fingerprinting.png}

\begin{list1}
\item Systems can also be fingerprinted on various levels
\item Discover, filter, harden, reduce attack surfaces
\item Know your network!
\end{list1}


\slide{Most firewalls include some detection today}

\hlkimage{5cm}{NSM_Phases-300x238.png}

Source: ANSM chapter 1: The Practice of Applied Network Security Monitoring
\begin{list2}
\item Vulnerability-Centric vs. Threat-Centric Defense
\item The NSM cycle: collection, detection, and analysis
\item Full Content Data, Session Data, Statistical Data, Packet String Data, and Alert Data
\item Security Onion is nice, but a bit over the top - quickly gets overloaded

\item Book referenced is: \emph{Network Security Monitoring
Applied Network Security Monitoring Collection, Detection, and Analysis}, 2014 Chris Sanders ISBN: 9780124172081
\end{list2}



\slide{Indicators of Compromise and Signatures}

\begin{quote}
An IOC is any piece of information that can be used to objectively describe a
network intrusion, expressed in a platform-independent manner. This could include a simple indicator such as the IP address of a command and control (C2) server or a complex set of behaviors that indicate that a mail server is being used as a malicious SMTP relay.

When an IOC is taken and used in a platform-specific language or format, such as a Snort Rule or a Bro-formatted file, it becomes part of a signature. A signature can contain one or more IOCs.
\end{quote}

Source: Applied Network Security Monitoring Collection, Detection, and Analysis, 2014 Chris Sanders


\begin{list2}
\item Network Security Monitoring (NSM) - monitoring networks for intrusions, and reacting to those
\item networkbased intrusion detection systems (NIDS)
\item host based intrusion detection systems (HIDS)
\item Example systems are Security Onion \link{https://securityonion.net/} or\\ SELKS \link{https://www.stamus-networks.com/open-source/}
\end{list2}



\slide{Network Sniffing for Security}


\begin{list1}
\item ANSM chapter 3:The Sensor Platform
\begin{list2}
\item Full Packet Capture (FPC) Data
\item Session Data
\item Statistical Data
\item Packet String (PSTR) Data
\item Log Data
\item Sensor Placement, designing etc.
\end{list2}
\end{list1}
\emph{Applied Network Security Monitoring Collection, Detection, and Analysis}, 2014 Chris Sanders ISBN: 9780124172081 - shortened ANSM


\slide{Collect Network Evidence from the network -- netflow}

\begin{list1}
\item Netflow is getting more important, more data share the same links
\item Detecting DoS/DDoS and problems is essential
\item Netflow sampling is vital information - 123Mbit, but what kind of traffic
\item Cisco standard NetFlow version 5 defines a flow as a unidirectional sequence of packets that all share the following 7 values:
\begin{list2}
\item Ingress interface (SNMP ifIndex), IP protocol, Source IP address and Destination IP address
\item Source port for UDP or TCP, 0 for other protocols, IP Type of Service
\item Destination port for UDP or TCP, type and code for ICMP, or 0 for other protocols
\end{list2}
\item Today we can use Netflow version 9 or IPFIX with more fields available
\end{list1}

Source: \\{\footnotesize
\link{https://en.wikipedia.org/wiki/NetFlow}
\link{https://en.wikipedia.org/wiki/IP_Flow_Information_Export}}


\slide{Architecture for packet capture}

\hlkimage{5cm}{network-horiz-onion.png}
Source: picture from \link{https://docs.securityonion.net/en/2.3/introduction.html}

\begin{list2}
\item Note the terminology North-South -- from the internet into the systems
\item East-West -- horizontal traffic inside the data center
\item See also from Security Onion \url{https://docs.securityonion.net/en/2.3/architecture.html#architecture}
\end{list2}



\slide{The Zeek Network Security Monitor}

\hlkimage{13cm}{zeek-overview.png}

The Zeek Network Security Monitor is not a single tool, more of a
powerful network analysis framework. Note: the project was renamed from Bro to Zeek in Oct 2018

{\small Source \url{https://www.zeek.org/}}


\slide{Suricata IDS/IPS/NSM}

\hlkimage{6cm}{suricata.png}

\begin{quote}
Suricata is a high performance Network IDS, IPS and Network Security Monitoring engine. Open Source and owned by a community run non-profit foundation, the Open Information Security Foundation (OISF). Suricata is developed by the OISF and its supporting vendors.
\end{quote}

 \link{http://suricata-ids.org/}
 \link{http://openinfosecfoundation.org}

\centerline{I often use Suricata and Zeek together}

Recommend the Emerging Threats ruleset, ET Open is free, ET Pro is a commercial license\\
\url{https://rules.emergingthreats.net/}



\slide{Basic port scanning}

Testing your TCP/IP networks, map them!

\begin{list2}
\item What is a port scan
\item Testing all values possible for port number from 0/1 to 65535
\item Goal is to identify open ports, listening and vulnerable services
\item Most often TCP and UDP scan
\item TCP scanning is more reliable than UDP scanning
\item TCP handshake must respond with SYN-ACK packets
\item UDP applications respond differently -- if they even respond, so probes with real requests may get response, no firewall they respond with ICMP on closed ports
\end{list2}

\slide{TCP three-way handshake}

\hlkimage{4cm}{images/tcp-three-way.pdf}

\begin{list2}
\item {\bfseries TCP SYN half-open} scans
\item in the old days systems would only log when a full TCP connection was setup
  -- so doing only half open it was a \emph{stealth}-scans
\item Today system and IDS intrusion detection can easily monitor for this
\item Sending a lot of SYN packets can create a Denial of Service -- {\bfseries SYN-flooding}
\end{list2}

\slide{Scope: select systems for testing}

\hlkimage{10cm}{overview-routing-customer-2015.png}

\begin{list2}
\item Routers in front of critical systems and networks - availability
\item Firewalls -- are traffic flows restricted
\item Mail servers -- open for relaying
\item Web servers -- remote code execution in web systems, data download
\end{list2}

\slide{Ping and port sweep}

\begin{list1}
\item Scans across the network are named sweeps
\item Ping sweeps using ICMP Ping probes
\item Port sweep trying to find a specific service, like port 80 web
\item Quite easy to see in network traffic:
\begin{list2}
\item Selecting two IP-addresses not in use
\item Should not see any traffic, but if it does, its being scanned
\item If traffic is received on both addresses, its a sweep -- if they are a bit apart it is even better, like 10.0.0.100 and 10.0.0.200
  \end{list2}

\vskip 2cm
Pro tip: a Great network intrusion detection engine (IDS), is Suricata \link{https://suricata.io}
\end{list1}

\slide{what is Nmap today}
\begin{quote}
Nmap ("Network Mapper") is a free and open source (license) utility for network discovery and security auditing.
\end{quote}

\begin{list1}
\item Initial release September 1997 -- almost 30 years ago
\item Today a package of programs for Windows, Mac, BSD, Linux, ... source
\item Flexible, powerful, and free! Includes other tools!
\item Check the changelog for new releases\\
\link{https://nmap.org/changelog.html}
\end{list1}

Bonus info: you can help Nmap by submitting fingerprints


\slide{Nmap port sweep for web servers}

\begin{alltt}\small
root@cornerstone:~#{\bfseries  nmap -p80,443 172.29.0.0/24}

Starting Nmap 6.47 ( http://nmap.org ) at 2015-02-05 07:31 CET
Nmap scan report for 172.29.0.1
Host is up (0.00016s latency).
PORT    STATE    SERVICE
{\color{darkgreen}80/tcp  open     http}
443/tcp filtered https
MAC Address: 00:50:56:C0:00:08 (VMware)

Nmap scan report for 172.29.0.138
Host is up (0.00012s latency).
PORT    STATE  SERVICE
{\color{darkgreen}80/tcp  open   http}
443/tcp closed https
MAC Address: 00:0C:29:46:22:FB (VMware)

\end{alltt}

\slide{Nmap port sweep for SNMP port 161/UDP}

\begin{alltt}\small
root@cornerstone:~#{\bfseries nmap -sU -p 161 172.29.0.0/24}
Starting Nmap 6.47 ( http://nmap.org ) at 2015-02-05 07:30 CET
Nmap scan report for 172.29.0.1
Host is up (0.00015s latency).
PORT    STATE         SERVICE
{\color{darkgreen}161/udp open|filtered snmp}
MAC Address: 00:50:56:C0:00:08 (VMware)

Nmap scan report for 172.29.0.138
Host is up (0.00011s latency).
PORT    STATE  SERVICE
{\bf{161/udp closed snmp}}
MAC Address: 00:0C:29:46:22:FB (VMware)
...
Nmap done: 256 IP addresses (5 hosts up) scanned in 2.18 seconds
\end{alltt}

\vskip 5mm
\centerline{More reliable to use Nmap script with probes like --script=snmp-info}

\slide{Nmap Advanced OS detection}
\begin{alltt}\footnotesize
root@cornerstone:~#{\bfseries nmap -A -p80,443 172.29.0.0/24}
Starting Nmap 6.47 ( http://nmap.org ) at 2015-02-05 07:37 CET
Nmap scan report for 172.29.0.1
Host is up (0.00027s latency).
PORT    STATE    SERVICE VERSION
80/tcp  open     http    Apache httpd 2.2.26 ((Unix) DAV/2 mod_ssl/2.2.26 OpenSSL/0.9.8zc)
|_http-title: Site doesn't have a title (text/html).
443/tcp filtered https
MAC Address: 00:50:56:C0:00:08 (VMware)
Device type: media device|general purpose|phone
Running: Apple iOS 6.X|4.X|5.X, Apple Mac OS X 10.7.X|10.9.X|10.8.X
OS details: Apple iOS 6.1.3, Apple Mac OS X 10.7.0 (Lion) - 10.9.2 (Mavericks)
or iOS 4.1 - 7.1 (Darwin 10.0.0 - 14.0.0), Apple Mac OS X 10.8 - 10.8.3 (Mountain Lion)
or iOS 5.1.1 - 6.1.5 (Darwin 12.0.0 - 13.0.0)
OS and Service detection performed.
Please report any incorrect results at http://nmap.org/submit/
\end{alltt}

\begin{list2}
\item Low-level way to identify operating systems, also try/use
  \verb+nmap -A+
\item Send probes and observe responses, lookup in table of known OS and responses
\item Techniques known since at least: \emph{ICMP Usage In Scanning} Version 3.0,
  Ofir Arkin, 2001 %\link{https://web.archive.org/web/20050210093427/http://www.sys-security.com/html/projects/icmp.html} % Original side er død
\end{list2}

\exercise{ex:wireshark-install}
\exercise{ex:nmap-pingsweep}
\exercise{ex:nmap-synscan}
\exercise{ex:nmap-os}
\exercise{ex:web-site-check}
\exercise{ex:zeekweb}


\slide{Conclusion -- Networks are exciting!}

\hlkimage{10cm}{IMG_20250309_bailey.jpg}


\begin{list2}
\item Many protocols work together to provide \emph{cat and dog pictures}
\item Lots of acronyms, but a relatively small set of protocols will get you started
\item Working with a small portable router and your laptop might teach you the basics
\end{list2}

%Pictured a small white and black dog playing with a toy porcupine
\myquestionspage






\end{document}
